\Needspace{9\baselineskip}\begin{center}\itshape \textit{Anwendung der vorstehenden Untersuchungen auf die trigonometrischen Funktionen.}\\\textit{Eine Methode, die Winkel zu bestimmen, denen die einzelnen Wurzeln \(\Omega\) entsprechen.}\end{center}


\Needspace{6\baselineskip}\Article{361}
Es bleibt uns noch, den Zusammenhang zwischen den Wurzeln \(\Omega\) und den trigonometrischen Funktionen der Winkel \(\frac Pn,\allowbreak \frac{2P}n,\allowbreak \frac{3P}n\ldots\frac{(n-1)P}n\) näher zu betrachten. Die Methode, die wir zum Auffinden der Wurzeln \(\Omega\) angegeben haben, lässt noch ungewiss (sofern nicht während der Arbeit die Sinustafeln in der oben beschriebenen Weise benutzt wurden, was jedoch weniger unmittelbar wäre), \textit{welche} Wurzeln den \textit{einzelnen} Winkeln entsprechen, \(\text{d.\,h.}\) welche Wurzel \(=\cos\frac Pn+i\sin\frac Pn\), welche \(=\cos\frac{2P}n+i\sin\frac{2P}n\) ist, usw. Diese Ungewissheit lässt sich leicht beheben, wenn man bedenkt, dass die Kosinus der Winkel \(\frac Pn,\allowbreak \frac{2P}n,\allowbreak \frac{3P}n\ldots\frac{(n-1)P}{2n}\) fortwährend abnehmen (sofern man auch die Vorzeichen berücksichtigt) und ihre Sinus alle positiv sind; die Winkel \(\frac{(n-1)P}n,\allowbreak \frac{(n-2)P}n,\allowbreak \frac{(n-3)P}n\ldots\frac{(n+1)P}{2n}\) dagegen haben der Reihe nach dieselben Kosinus wie jene, aber negative Sinus, deren absolute Beträge den Sinus jener Winkel gleich sind. Daher entsprechen von den Wurzeln \(\Omega\) diejenigen beiden mit den größten reellen Teilen (die einander gleich sind) den Winkeln \(\frac Pn,\allowbreak \frac{(n-1)P}n\): dem ersten die Wurzel, in der die imaginäre Größe \(i\) mit einer positiven Größe multipliziert ist, dem zweiten diejenige, in der \(i\) mit einer negativen Größe multipliziert ist. Von den \(n-3\) übrigen Wurzeln entsprechen wiederum diejenigen mit den größten reellen Teilen den Winkeln \(\frac{2P}n,\allowbreak \frac{(n-2)P}n\), und so weiter. — Sobald die Wurzel, der der Winkel \(\frac Pn\) entspricht, erkannt ist, lassen sich diejenigen, die den übrigen Winkeln entsprechen, auch dadurch unterscheiden, dass, wenn jene \(=[\lambda]\) gesetzt wird, den Winkeln \(\frac{2P}n,\allowbreak \frac{3P}n,\allowbreak \frac{4P}n\), usw., offenbar die Wurzeln \([2\lambda],\allowbreak [3\lambda],\allowbreak [4\lambda]\), usw., entsprechen. So sieht man im Beispiel von Artikel 353 sofort, dass dem Winkel \(\frac1{19}P\) keine andere Wurzel als \([11]\) und dem Winkel \(\frac{18}{19}P\) die Wurzel \([8]\) entsprechen kann; ebenso entsprechen den Winkeln \(\frac2{19}P,\allowbreak \frac{17}{19}P,\allowbreak \frac3{19}P,\allowbreak \frac{16}{19}P\), usw., die Wurzeln \([3],\allowbreak [16],\allowbreak [14],\allowbreak [5]\), usw. Im Beispiel von Artikel 354 entspricht dem Winkel \(\frac1{17}P\) offenbar die Wurzel \([1]\), dem Winkel \(\frac2{17}P\) die Wurzel \([2]\), usw. Auf diese Weise sind also die Kosinus und Sinus der Winkel \(\frac Pn,\allowbreak \frac{2P}n\), usw., vollständig bestimmt.
\SourceLeaf{460}

\clearpage
\begingroup\abovedisplayskip=6pt plus2pt minus1pt \belowdisplayskip=6pt plus2pt minus1pt \abovedisplayshortskip=0pt plus1pt \belowdisplayshortskip=3pt plus1pt minus1pt
\Needspace{9\baselineskip}\begin{center}\itshape \textit{Tangens, Kotangens, Sekans und Kosekans werden aus Sinus und Kosinus ohne Division abgeleitet.}\end{center}


\Needspace{6\baselineskip}\Article{362}
Was die übrigen trigonometrischen Funktionen dieser Winkel angeht, so ließen sie sich aus den entsprechenden Kosinus und Sinus mit den allgemein bekannten Methoden leicht ableiten: nämlich die Sekans- und Tangenswerte, indem man die Einheit und die Sinus durch die Kosinus dividiert, sowie die Kosekans- und Kotangenswerte, indem man die Einheit und die Kosinus durch die Sinus dividiert. Bequemer erhält man jedoch dasselbe meist mit den folgenden Formeln, ohne Divisionen, allein durch Additionen. Es sei \(\omega\) irgendeiner der Winkel \(\frac Pn,\allowbreak \frac{2P}n\ldots\frac{(n-1)P}n\), und es sei \(\cos\omega+i\sin\omega=R\), sodass \(R\) eine der Wurzeln \(\Omega\) ist; dabei gilt \[\begin{aligned}\cos\omega&=\tfrac12(R+\tfrac1R)=\frac{1+RR}{2R},\\\sin\omega&=\tfrac1{2i}(R-\tfrac1R)=\frac{i(1-RR)}{2R}\end{aligned}\] Daraus folgt \[\begin{gathered}\sec\omega=\frac{2R}{1+RR}.\quad\operatorname{tang}\omega=\frac{i(1-RR)}{1+RR},\\\operatorname{cosec}\omega=\frac{2Ri}{RR-1},\quad\operatorname{cotang}\omega=\frac{i(RR+1)}{RR-1}\end{gathered}\] Nun wollen wir zeigen, wie man die Zähler dieser vier Brüche so umformen kann, dass sie durch ihre Nenner teilbar werden.

I. Wegen \(R=R^{n+1}=R^{2n+1}\) wird \(2R=R+R^{2n+1}\), welcher Ausdruck offenbar durch \(1+RR\) teilbar ist, da \(n\) ungerade ist. Daraus folgt \[\sec\omega=R-R^3+R^5-R^7\ldots+R^{2n-1}\] und daher (da wegen \(\sin\omega=-\sin(2n-1)\omega,\allowbreak \ \sin3\omega=-\sin(2n-3)\omega\), usw., offenbar \(\sin\omega-\sin3\omega+\sin5\omega\ldots+\sin(2n-1)\omega=0\) gilt) \[\sec\omega=\cos\omega-\cos3\omega+\cos5\omega\ldots+\cos(2n-1)\omega\] oder schließlich (da \(\cos\omega=\cos(2n-1)\omega,\allowbreak \ \cos3\omega=\cos(2n-3)\omega\), usw.), \[=2(\cos\omega-\cos3\omega+\cos5\omega\ldots\mp\cos(n-2)\omega)\pm\cos n\omega\] wobei das obere oder untere Vorzeichen gilt, je nachdem \(n\) die Form \(4k+1\) oder \(4k+3\) hat. Offenbar lässt sich diese Formel auch so darstellen: \[\sec\omega=\pm(1-2\cos2\omega+2\cos4\omega\ldots\pm2\cos(n-1)\omega)\]

II. Setzt man ebenso \(1-R^{2n+2}\) für \(1-RR\) ein, so ergibt sich
\SourceLeaf{461}
\[\operatorname{tang}\omega=i(1-RR+R^4-R^6\ldots-R^{2n})\] oder (da \(1-R^{2n}=0,\allowbreak \ RR-R^{2n-2}=2i\sin2\omega,\allowbreak \ R^4-R^{2n-4}=2i\sin4\omega\), usw.), \[\operatorname{tang}\omega=2(\sin2\omega-\sin4\omega+\sin6\omega\ldots\mp\sin(n-1)\omega)\]

III. Da \(1+RR+R^4\ldots+R^{2n-2}=0\) gilt, wird \[\begin{aligned}n&=n-1-RR-R^4\ldots R^{2n-2}\\&=(1-1)+(1-RR)+(1-R^4)\ldots+(1-R^{2n-2})\end{aligned}\] und jeder einzelne Teil dieser Summe ist durch \(1-RR\) teilbar. Daraus folgt \[\begin{aligned}\frac n{1-RR}&=1+(1+RR)+(1+RR+R^4)\ldots\\&\quad+(1+RR+R^4\ldots+R^{2n-4})\\&=(n-1)+(n-2)RR+(n-3)R^4\ldots+R^{2n-4}\end{aligned}\] sodass man durch Multiplikation mit \(2\), Subtraktion von \[0=(n-1)(1+RR+R^4\ldots+R^{2n-2})\] und erneute Multiplikation mit \(R\) \[\begin{aligned}\frac{2nR}{1-RR}&=(n-1)R+(n-3)R^3+(n-5)R^5\ldots\\&\quad-(n-3)R^{2n-3}-(n-1)R^{2n-1}\end{aligned}\] erhält, woraus sich sofort ergibt: \[\begin{aligned}\operatorname{cosec}\omega&=\frac1n\bigl((n-1)\sin\omega+(n-3)\sin3\omega\ldots\\&\qquad\qquad-(n-1)\sin(2n-1)\omega\bigr)\\&=\frac2n\bigl((n-1)\sin\omega+(n-3)\sin3\omega\\&\qquad\qquad+\text{usw.}+2\sin(n-2)\omega\bigr)\end{aligned}\] Diese Formel lässt sich auch so darstellen: \[\operatorname{cosec}\omega=-\frac2n(2\sin2\omega+4\sin4\omega+6\sin6\omega\ldots+(n-1)\sin(n-1)\omega)\]

IV. Durch Multiplikation des oben angegebenen Wertes von \(\frac n{1-RR}\) mit \(1+RR\) und Subtraktion von \[0=(n-1)(1+RR+R^4\ldots+R^{2n-2})\] ergibt sich \[\frac{n(1+RR)}{1-RR}=(n-2)RR+(n-4)R^4+(n-6)R^6\ldots-(n-2)R^{2n-2}\]
\SourceLeaf{462}
woraus sofort folgt: \[\begin{aligned}\operatorname{cotang}\omega&=\frac1n\bigl((n-2)\sin2\omega+(n-4)\sin4\omega\\&\qquad+(n-6)\sin6\omega\ldots-(n-2)\sin(n-2)\omega\\&=\frac2n((n-2)\sin2\omega+(n-4)\sin4\omega\ldots\\&\qquad+3\sin(n-3)\omega+\sin(n-1)\omega)\end{aligned}\] Diese Formel lässt sich auch folgendermaßen darstellen: \[\operatorname{cotang}\omega=-\frac2n(\sin\omega+3\sin3\omega\ldots+(n-2)\sin(n-2)\omega)\]\endgroup


